\documentclass[../../main.tex]{subfiles}

\begin{document}

\section{良序集}\label{section:ch01-良序集}

正整数集 $\mathbb{Z}_{+}$ 有一个有用的性质：每一个非空子集有一个最小元。将它加以推广就得到良序集概念。

\begin{definition}[良序集]\label{definition:ch01-良序集}
具有全序关系 $<$ 的一个集合 $A$ 称为良序的（well-ordered），如果 $A$ 的任意非空子集有一个最小元。
\end{definition}

\begin{example}\label{example:ch01-良序例子1}
在字典序下考虑集合 $\{1, 2\} \times \mathbb{Z}_{+}$。它可以表示成一个无穷序列紧跟着另一个无穷序列：
\begin{align*}
a_{1}, a_{2}, a_{3}, \cdots; b_{1}, b_{2}, b_{3}, \cdots,
\end{align*}
其中每一个元素都比它右边的每一个元素要小。不难看出，这个全序集的每一个非空子集 $C$ 都有一个最小元。如果 $C$ 包含任何一个元素 $a_{n}$，这个最小元便是 $C$ 与序列 $a_{1}, a_{2}, \cdots$ 的交的最小元；而如果 $C$ 不包含 $a_{n}$，那么 $C$ 是序列 $b_{1}, b_{2}, \cdots$ 的一个子集，因此有一个最小元。
\end{example}

\par
\begin{figure}[htbp]
\centering
\includegraphics[width=0.6\textwidth]{ch01-fig10-1}
\caption{$\mathbb{Z}_{+} \times \mathbb{Z}_{+}$ 字典序下的最小元}\label{figure:ch01-良序最小元}
\end{figure}

\begin{example}\label{example:ch01-良序例子2}
在字典序下考虑集合 $\mathbb{Z}_{+} \times \mathbb{Z}_{+}$。它可以表示成一个无穷序列的无穷序列。我们断定它是一个良序集。设 $X$ 是 $\mathbb{Z}_{+} \times \mathbb{Z}_{+}$ 的一个非空子集，$A$ 是 $\mathbb{Z}_{+}$ 中由所有 $X$ 的元素的第一个坐标组成的子集。$A$ 有最小元，设为 $a_{0}$，则族
\begin{align*}
\{b \mid a_{0} \times b \in X\}
\end{align*}
是 $\mathbb{Z}_{+}$ 的一个非空子集。设 $b_{0}$ 为其最小元。根据字典序的定义，$a_{0} \times b_{0}$ 就是 $X$ 的最小元（参见\cref{figure:ch01-良序最小元}）。

\end{example}

\begin{example}\label{example:ch01-非良序例子}
整数集在通常全序关系下不是良序集。这是因为由负整数组成的子集没有最小元。区间 $0 \leqslant x \leqslant 1$ 中的全体实数也不是良序集，这是因为由 $0 < x < 1$ 中所有 $x$ 组成的子集没有最小元（尽管它有一个下确界）。
\end{example}

有几种构造良序集的方法，下面是其中的两种：
\begin{enumerate}[(1)]
\item 如果 $A$ 是一个良序集，则 $A$ 的任意子集在限制全序关系下是一个良序集.
\item 如果 $A$ 和 $B$ 都是良序集，则 $A \times B$ 在字典序下是一个良序集.
\end{enumerate}
(1) 的证明是显然的，(2) 的证明可以按\cref{example:ch01-良序例子2}的方式给出。

由此可得，$\mathbb{Z}_{+} \times (\mathbb{Z}_{+} \times \mathbb{Z}_{+})$ 在字典序下是良序的，它可以表示成一个无穷序列的无穷序列的无穷序列。类似地，$(\mathbb{Z}_{+})^{4}$ 在字典序下是良序的，等等。如果你试图把它推广到无穷多个 $\mathbb{Z}_{+}$ 的积，那就要遇到麻烦了。我们将要简单地讨论这种情况。

给定一个没有全序关系的集合 $A$，人们自然会问，在 $A$ 上是否存在一个全序关系，使其成为良序集呢？如果 $A$ 是有限集，那么任意一一对应
\begin{align*}
f: A \longrightarrow \{1, \cdots, n\}
\end{align*}
就能够定义 $A$ 上的一个全序关系，使 $A$ 与全序集 $\{1, \cdots, n\}$ 的序型相同。事实上，有限集上的任意全序关系都可以用上述方法得到。

\begin{theorem}[有限全序集的序型]\label{theorem:ch01-有限良序}
任意非空有限全序集具有 $\mathbb{Z}_{+}$ 的一个截 $\{1, \cdots, n\}$ 的序型，因而必定是良序集。
\end{theorem}

\begin{proof}
这个定理已在第\cref{section:ch01-有限集}节中作为一个习题给出，现在予以证明。首先证明任意有限全序集 $A$ 有一个最大元。如果 $A$ 只有一个元素，这是显然的。假定结论对于具有 $n-1$ 个元素的集合成立，设 $A$ 有 $n$ 个元素并且 $a_{0} \in A$。那么 $A - \{a_{0}\}$ 有一个最大元 $a_{1}$，而 $\{a_{0}, a_{1}\}$ 的最大者就是 $A$ 的最大元。

其次证明，对于某一个 $n$，存在一个 $A$ 与 $\{1, \cdots, n\}$ 之间的保序一一对应。如果 $A$ 有一个元素，这是显然的。假定结论对于有 $n-1$ 个元素的集合成立。设 $b$ 为 $A$ 的最大元。根据归纳假定，存在一个保序一一对应
\begin{align*}
f': A - \{b\} \longrightarrow \{1, \cdots, n-1\}.
\end{align*}
再由下式
\begin{align*}
f(x) = f'(x), x \neq b,
\end{align*}
\begin{align*}
f(b) = n
\end{align*}
定义一个保序一一对应 $f: A \longrightarrow \{1, \cdots, n\}$。
\end{proof}

因此，一个有限全序集只有一种可能的序型。对于无限集，情形就完全不同了。良序集
\begin{align*}
\mathbb{Z}_{+},
\end{align*}
\begin{align*}
\{1, \cdots, n\} \times \mathbb{Z}_{+},
\end{align*}
\begin{align*}
\mathbb{Z}_{+} \times \mathbb{Z}_{+},
\end{align*}
\begin{align*}
\mathbb{Z}_{+} \times (\mathbb{Z}_{+} \times \mathbb{Z}_{+})
\end{align*}
都是可数无限集，而它们有不同的序型。

我们已经给出的所有良序集的例子都是可数集。人们自然会问，能不能找到一个不可数的良序集呢？

大家熟悉的不可数集是可数无穷多个 $\mathbb{Z}_{+}$ 的积
\begin{align*}
X = \mathbb{Z}_{+} \times \mathbb{Z}_{+} \times \cdots = (\mathbb{Z}_{+})^{\omega}.
\end{align*}
可以用一种自然的方式在这个集合中引进广义的字典序：如果对于某一个 $n \geqslant 1$，
\begin{align*}
\text{当 } i < n \text{ 时, } a_{i} = b_{i}; \text{ 并且 } a_{n} < b_{n},
\end{align*}
则定义
\begin{align*}
(a_{1}, a_{2}, \cdots) < (b_{1}, b_{2}, \cdots).
\end{align*}
事实上，这是集合 $X$ 上的一个全序关系。可惜它不是一个良序。考虑 $X$ 中所有形如
\begin{align*}
\boldsymbol{x} = (1, \cdots, 1, 2, 1, 1, \cdots)
\end{align*}
的元素 $\boldsymbol{x}$ 的集合 $A$，其中 $\boldsymbol{x}$ 只有一个坐标等于 $2$，其余的都是 $1$。显然 $A$ 没有最小元。

我们已经看到，字典序至少没有给出集合 $(\mathbb{Z}_{+})^{\omega}$ 的一个良序。那么在这个集合上是不是有别的全序关系，使它成为良序集呢？我们虽然没有一个在 $(\mathbb{Z}_{+})^{\omega}$ 上作出特定良序的方法，但是有一个著名的定理，可确认这样的良序是存在的。

\begin{theorem}[良序定理]\label{theorem:ch01-良序性质定理}
若 $A$ 是一个集合，则存在 $A$ 上的一个全序关系，使 $A$ 成为一个良序集。
\end{theorem}

这个定理是 1904 年由 Zermelo 证明的，曾震动了数学界，并对其证明的正确性展开大量辩论。对于任意不可数集，没有一个将其良序化的构造程序就引起了许多怀疑。在进一步分析这个证明时，我们发现只有一处可能有点问题，就是在构造的过程中包含着无穷多次选择，也就是说其做法包含了\cref{axiom:ch01-选择公理}。

一些数学家曾经拒绝把\cref{axiom:ch01-选择公理}作为一个结果，多少年来每涉及一个新定理就要提出这个问题：它的证明是否涉及\cref{axiom:ch01-选择公理}？因为一个定理的证明中如果使用了\cref{axiom:ch01-选择公理}，便会被认为基础不稳固。一般说来，时下的数学家大都已经不再有这样的烦恼了。他们把\cref{axiom:ch01-选择公理}作为集合论中的一个合适的假定，良序定理也就随之得以被承认。

\cref{axiom:ch01-选择公理}蕴涵良序定理的证明过于冗长（虽然不是特别困难），而且主要是逻辑学家的兴趣所在，我们略去其证明。如果你有兴趣的话，可以在本章最后的附加习题中找到证明的纲要。本书承认良序定理，读者也可以将其视作另外的一个公理！

我们只是偶尔需要良序定理的全部“功效”。一般情况下，我们将仅用到下面这个较弱的结果：

\begin{corollary}[存在不可数良序集]\label{corollary:ch01-不可数良序}
存在一个不可数的良序集。
\end{corollary}

我们将使用这个结果来构造一个很有用的良序集。

\begin{definition}[截]\label{definition:ch01-截}
设 $X$ 是一个全序集。给定 $\alpha \in X$，用 $S_{\alpha}$ 表示集合
\begin{align*}
S_{\alpha} = \{x \mid x \in X \text{ 并且 } x < \alpha\}.
\end{align*}
称之为 $X$ 在 $\alpha$ 处的截（section）。
\end{definition}

\begin{lemma}[极小不可数良序集]\label{lemma:ch01-极小不可数}
存在一个以 $\Omega$ 为最大元的不可数良序集 $A$，$A$ 在 $\Omega$ 处的截 $S_{\Omega}$ 是一个不可数集，而 $A$ 的每一个其余的截都是可数集。
\end{lemma}

\begin{proof}
假定 $B$ 是一个不可数的良序集。设 $C$ 为 $\{1, 2\} \times B$ 在字典序下的不可数的良序集，则 $C$ 的某一个截不可数。（事实上，$C$ 在每一个形如 $2 \times b$ 的元素处的截都是不可数的。）记 $\Omega$ 为使得 $C$ 在 $\Omega$ 处的截为不可数集的最小元。取 $A$ 为这个截加上 $\Omega$ 所组成的集合。
\end{proof}

值得注意的是，$S_{\Omega}$ 是一个不可数良序集，并且它的每一个截都是一个可数集。事实上，它的序型由此而唯一确定。我们称之为极小不可数良序集。进而，我们将用记号 $\bar{S}_{\Omega}$ 来表示不可数良序集 $A = S_{\Omega} \cup \{\Omega\}$（理由后述）。

就我们的目的而言，$S_{\Omega}$ 的最有用的性质是以下定理：

\begin{theorem}[可数子集有上界]\label{theorem:ch01-可数上界}
如果 $A$ 是 $S_{\Omega}$ 的一个可数子集，则 $A$ 在 $S_{\Omega}$ 中有上界。
\end{theorem}

\begin{proof}
设 $A$ 为 $S_{\Omega}$ 的一个可数子集。对于每一个 $a \in A$，截 $S_{a}$ 是可数的。因此，并 $B = \bigcup_{a \in A} S_{a}$ 也是可数的。因为 $S_{\Omega}$ 不可数，所以集合 $B$ 不会等于 $S_{\Omega}$。设 $x$ 为 $S_{\Omega}$ 中不属于 $B$ 的点，则 $x$ 为 $A$ 的一个上界。因为，如果对于 $A$ 中某一个 $a$ 有 $x < a$ 成立，那么 $x$ 属于 $S_{a}$，从而属于 $B$，这与 $x$ 的选取矛盾。
\end{proof}

\subsection*{习题}

\begin{exercise}\label{exercise:ch01-良序集-1}
证明：每一个良序集都有上确界性质。
\end{exercise}

\begin{exercise}\label{exercise:ch01-良序集-2}
(a) 证明：在良序集中，每一个不是最大元（如果存在的话）的元素有一个紧接后元。

(b) 作出一个集合，它的每一个元素有一个紧接后元，但这个集合不是良序集。
\end{exercise}

\begin{exercise}\label{exercise:ch01-良序集-3}
集合 $\{1, 2\} \times \mathbb{Z}_{+}$ 与 $\mathbb{Z}_{+} \times \{1, 2\}$ 在字典序下都是良序集，它们有相同的序型吗？
\end{exercise}

\begin{exercise}\label{exercise:ch01-良序集-4}
(a) 用 $\mathbb{Z}_{-}$ 表示常用全序关系下的负整数集。证明：一个全序集 $A$ 不是良序集的充分必要条件是它包含一个与 $\mathbb{Z}_{-}$ 有相同序型的子集。

(b) 证明：若 $A$ 是一个全序集并且其每一个可数子集都是良序集，则 $A$ 是一个良序集。
\end{exercise}

\begin{exercise}\label{exercise:ch01-良序集-5}
证明良序定理蕴涵着\cref{axiom:ch01-选择公理}。
\end{exercise}

\begin{exercise}\label{exercise:ch01-良序集-6}
设 $S_{\Omega}$ 为极小不可数良序集。
\begin{enumerate}[(a)]
\item 证明：$S_{\Omega}$ 没有最大元.
\item 证明：对于每一个 $\alpha \in S_{\Omega}$，集合 $\{x \mid \alpha < x\}$ 是不可数的.
\item 设 $X_{0}$ 为 $S_{\Omega}$ 的一个子集，它由所有没有紧接前元的元素 $x$ 组成。证明 $X_{0}$ 是不可数的.
\end{enumerate}
\end{exercise}

\begin{exercise}\label{exercise:ch01-良序集-7}
设 $J$ 是一个良序集。$J$ 的子集 $J_{0}$ 称为一个归纳集（inductive set），如果对于任意 $\alpha \in J$，
\begin{align*}
(S_{\alpha} \subset J_{0}) \Longrightarrow \alpha \in J_{0}.
\end{align*}

\begin{theorem}[超限归纳原理]\label{theorem:ch01-超限归纳}
若 $J$ 是一个良序集，并且 $J_{0}$ 是 $J$ 的一个归纳子集，则 $J_{0} = J$。
\end{theorem}
\end{exercise}

\begin{exercise}\label{exercise:ch01-良序集-8}
(a) 设 $A_{1}$ 和 $A_{2}$ 是两个无交的集合，并且分别对于序 $<_{1}$ 和序 $<_{2}$ 而言是良序集。在 $A_{1} \cup A_{2}$ 上定义一个全序关系“$<$”：如果有 $a, b \in A_{1}$ 和 $a <_{1} b$；或者有 $a, b \in A_{2}$ 和 $a <_{2} b$；或者有 $a \in A_{1}$ 和 $b \in A_{2}$ 时；定义 $a < b$。证明这是一个良序。

(b) 把 (a) 推广到无交良序集的任意一个族上去，其指标集为良序集。
\end{exercise}

\begin{exercise}\label{exercise:ch01-良序集-9}
考虑 $(\mathbb{Z}_{+})^{\omega}$ 的子集 $A$，这里 $A$ 由尾部全是 $1$ 的正整数序列 $\boldsymbol{x} = (x_{1}, x_{2}, \cdots)$ 所构成。在 $A$ 上定义一个序：如果 $x_{n} < y_{n}$ 那么 $\boldsymbol{x} < \boldsymbol{y}$ 并且当 $i > n$ 时，$x_{i} = y_{i}$。我们将此序称之为 $A$ 上的“反字典序”。
\begin{enumerate}[(a)]
\item 证明：对于每一个 $n$，存在 $A$ 的一个截与在字典序下的 $(\mathbb{Z}_{+})^{n}$ 有相同的序型.
\item 证明 $A$ 是良序的.
\end{enumerate}
\end{exercise}

\begin{exercise}\label{exercise:ch01-良序集-10}
定理 设 $J$ 与 $C$ 是两个良序集。假定不存在从 $J$ 的一个截到 $C$ 上的满射。则对于每一个 $x \in J$，存在唯一的一个函数 $h: J \longrightarrow C$ 满足
\begin{align}
h(x) = [C - h(S_{x})] \text{ 的最小元},\label{align:ch01-超限递归}
\end{align}
其中 $S_{x}$ 是 $J$ 在 $x$ 处的截。

证明：
\begin{enumerate}[(a)]
\item 如果 $h$ 和 $k$ 将 $J$ 的截或 $J$ 映入 $C$ 中，并且对于它们各自定义域中的所有 $x$ 满足\eqref{align:ch01-超限递归}，证明对于定义域中公共元 $x$ 有 $h(x) = k(x)$.
\item 如果存在满足\eqref{align:ch01-超限递归}的函数 $h: S_{\alpha} \longrightarrow C$，证明存在满足\eqref{align:ch01-超限递归}的函数 $k: S_{\alpha} \cup \{\alpha\} \longrightarrow C$.
\item 如果 $K \subset J$ 并且对于所有 $\alpha \in K$，存在满足\eqref{align:ch01-超限递归}的函数 $h_{\alpha}: S_{\alpha} \longrightarrow C$，证明存在满足\eqref{align:ch01-超限递归}的函数
\begin{align*}
k: \bigcup_{\alpha \in K} S_{\alpha} \longrightarrow C.
\end{align*}
\item 用超限归纳法证明，对于每一个 $\beta \in J$，存在满足\eqref{align:ch01-超限递归}的函数 $h_{\beta}: S_{\beta} \longrightarrow C$. [提示：若 $\beta$ 有紧接前元 $\alpha$，则 $S_{\beta} = S_{\alpha} \cup \{\alpha\}$。否则，$S_{\beta} = \bigcup_{\alpha < \beta} S_{\alpha}$。]
\item 证明本定理.
\end{enumerate}
\end{exercise}

\begin{exercise}\label{exercise:ch01-良序集-11}
设 $A$ 和 $B$ 是两个集合。用良序定理证明：或者它们有相同的基数，或者一个集合的基数大于另一集合的基数。[提示：若满射 $f: A \longrightarrow B$ 不存在，应用前面的习题。]
\end{exercise}

\end{document}